\documentclass{article}
\usepackage{amsmath,amssymb}
\usepackage{xurl}
\usepackage[hidelinks]{hyperref}
% Shared commands for notes in docs/paper-gaps/.

\DeclareUnicodeCharacter{2080}{\ifmmode{}_{0}\else\textsubscript{0}\fi}
\DeclareUnicodeCharacter{2081}{\ifmmode{}_{1}\else\textsubscript{1}\fi}
\DeclareUnicodeCharacter{2082}{\ifmmode{}_{2}\else\textsubscript{2}\fi}
\DeclareUnicodeCharacter{2099}{\ifmmode{}_{n}\else\textsubscript{n}\fi}

\newcommand{\Lean}{\textsc{Lean}}
\ExplSyntaxOn
\NewDocumentCommand{\leanid}{m}{
  \ifmmode
    \textsf{\detokenize{#1}}
  \else
    \group_begin:
    \urlstyle{sf}
    \tl_set:Nn \l_tmpa_tl {#1}
    \tl_replace_all:Nnn \l_tmpa_tl {\_} {_}
    \exp_args:NV \path \l_tmpa_tl
    \group_end:
  \fi
}
\ExplSyntaxOff
\providecommand{\tr}{\operatorname{tr}}
\newcommand{\tnleanIssueBase}{https://github.com/LionSR/TNLean/issues/}
\newcommand{\tnleanPrBase}{https://github.com/LionSR/TNLean/pull/}
\newcommand{\tnleanDocBase}{https://sirui-lu.com/TNLean/docs/}
\newcommand{\ghissue}[1]{\href{\tnleanIssueBase#1}{GitHub issue \##1}}
\newcommand{\ghpr}[1]{\href{\tnleanPrBase#1}{PR \##1}}
\newcommand{\doclink}[1]{\href{\tnleanDocBase#1.html}{\path{#1}}}

% Dirac notation and the identity operator, as in blueprint/src/macros/common.tex.
% \providecommand keeps note-local definitions authoritative.
\usepackage{dsfont}
\providecommand{\ket}[1]{|#1\rangle}
\providecommand{\bra}[1]{\langle #1|}
\providecommand{\braket}[2]{\langle #1 | #2 \rangle}
\providecommand{\ketbra}[2]{|#1\rangle\!\langle #2|}
\providecommand{\Id}{\mathds{1}}
\providecommand{\one}{\mathds{1}}
\providecommand{\C}{\mathbb{C}}
\providecommand{\MN}[1]{M_{#1}(\C)}
\providecommand{\MD}{M_D(\C)}

% External referencing for paper sources.  The build helper writes sanitized
% auxiliary files under build/paper-aux/ and excludes duplicated source labels.
\usepackage{xr}

% Import a generated paper auxiliary file with a caller-supplied label prefix.
% The candidate paths support builds from docs/paper-gaps/, the repository root,
% and build/paper-aux/ itself.
\newcommand{\importpaperaux}[2]{%
  \IfFileExists{../../build/paper-aux/#2.aux}{%
    \externaldocument[#1]{../../build/paper-aux/#2}%
  }{%
    \IfFileExists{build/paper-aux/#2.aux}{%
      \externaldocument[#1]{build/paper-aux/#2}%
    }{%
      \IfFileExists{#2.aux}{%
        \externaldocument[#1]{#2}%
      }{}%
    }%
  }%
}

% General external reference macros
\newcommand{\paperref}[2]{\ref{#1-#2}}
\newcommand{\papereqref}[2]{(\ref{#1-#2})}

% CPSV16 source (1606.00608 / MPDO-22-12-17-2.tex) external document and shortcuts
\importpaperaux{cpsv16-}{1606.00608/MPDO-22-12-17-2-tnlean}
\newcommand{\cpsvref}[1]{\paperref{cpsv16}{#1}}
\newcommand{\cpsveqref}[1]{\papereqref{cpsv16}{#1}}

% Verdict marker: \gapnote{<kind>}{<status>}. Kind is the policy
% classification (clarification, local-correction, scope-restriction,
% unfaithful, false-source, open-gap); status is open, wip (elimination
% underway), resolved, or historical. Machine-read by the site index;
% prints a verdict line.
\newcommand{\gapnote}[2]{\par\noindent\textbf{Verdict:} #1 (#2).\par}

\importpaperaux{fbc25-}{2502.20257/main-tnlean}
\title{CZX: Normalization and the Physical Z Gates}
\author{}
\date{2026-09-04}
\begin{document}
\maketitle
\gapnote{false-source}{open}

\section{The 2017 contraction convention and its obstruction}
Section II, ``Basics in Matrix Product Vectors'', of
\cite{Cirac2017MPU} defines the periodic contraction as an ordinary trace,
with no length-dependent prefactor. Section III, ``Structure of MPU'',
requires the resulting operator to be unitary and separately normalizes
its transfer matrix by $1/q$, where $q$ is the local physical dimension
($q=2$ here). That transfer normalization does not rescale the physical
operator.\footnote{Source traceability: \path{Papers/1703.09188/paper_v2.tex},
    lines 226--246 and 318--350. The declared auxiliary-label build for this
    archived source currently fails at the title's malformed accent command;
    the printed section and example designations are used here instead.}

Write $u,d,\ell,r\in\{0,1\}$ for upper, lower, left and right indices,
$X_{ud}=\delta_{u,1-d}$, and $H_{\ell r}=2^{-1/2}(-1)^{\ell r}$.
The example headed ``CZX'' in the time-reversal symmetry subsection of
\cite{Cirac2017MPU} prints\footnote{Source traceability:
    \path{Papers/1703.09188/paper_v2.tex}, lines 1914--1930.}
\begin{align}
    P^{ud}_{\ell r} & =X_{ud}\delta_{u,\ell}H_{\ell r}.
    \label{eq:czx_printed_tensor}
\end{align}
For binary output and input configurations $a,b$ of length $n>0$, with
cyclic indices $a_n=a_0$, the copy deltas force the unique virtual path
$\ell_j=a_j$. Consequently
\begin{align}
    \operatorname{MPO}_n(P)_{a,b}
     & =2^{-n/2}\prod_{j=0}^{n-1}\delta_{a_j,1-b_j}
    (-1)^{\sum_{j=0}^{n-1}a_ja_{j+1}}.
    \label{eq:czx_periodic_contraction}
\end{align}
Every row and column has one nonzero entry, of modulus $2^{-n/2}$, so
$\operatorname{MPO}_n(P)\operatorname{MPO}_n(P)^\dagger=2^{-n}I$,
not $I$. At $n=2$ the row indexed by $00$ has only the entry
$\operatorname{MPO}_2(P)_{00,11}=1/2$; its squared norm is $1/4$.
Blocking $L>0$ sites replaces $n$ by $Ln$ and cannot remove this obstruction.
This is a nondegenerate normalization error, not a zero-sector convention.

\section{What the 2025 diagram actually says}
Equations~\papereqref{fbc25}{eq:CZ} and
\papereqref{fbc25}{eq:deltas} of \cite{FrancoRubio2025MPUGauging}
define the blue two-leg tensor as $W=\sqrt{2}H$, with entries
$W_{\ell r}=(-1)^{\ell r}$; the wire junctions are ordinary copy deltas.
The CZX tensor in Equation~\papereqref{fbc25}{eq:MPU_CZX} of
\cite{FrancoRubio2025MPUGauging} has two such blue tensors, two lower
physical $X$ gates and two upper physical $Z$ gates. None may be
omitted.\footnote{Source traceability: \path{Papers/2502.20257/main.tex},
    lines 3795--3851 and 4503--4547.}

For comparison only, let $A=\sqrt{2}P$. Reading the two physical legs
left to right, let their outputs be $a,b$ and inputs $c,d$.
Writing $P_2$ and $A_2$ for the ordinary two-site blocks of $P$ and $A$,
the latter has entries
\begin{align}
    (A_2)^{ab,cd}_{\ell r}
     & =\delta_{a,1-c}\delta_{b,1-d}\delta_{\ell,a}(-1)^{ab+br}.
    \label{eq:czx_repaired_block}
\end{align}
Since $Z_{aa}=(-1)^a$, the displayed 2025 tensor is instead
\begin{align}
    B^{ab,cd}_{\ell r}
     & =\delta_{a,1-c}\delta_{b,1-d}\delta_{\ell,a}(-1)^{a+b+ab+br}
    =(-1)^{a+b}(A_2)^{ab,cd}_{\ell r}.
    \label{eq:czx_decorated_block}
\end{align}
Thus $B=(Z\otimes Z)A_2$ on the output physical leg, equivalently the
block of the one-site tensor with entries $(-1)^u A^{ud}_{\ell r}$.
It is neither $P_2$ nor $A_2$ in these coordinates. For example, at
$(a,b,c,d,\ell,r)=(0,1,1,0,0,0)$, their entries are respectively
$B=-1$, $P_2=1/2$, $A_2=1$.
On a ring of $m$ blocked sites the additional factor is
$(-1)^{\sum_{j=0}^{2m-1}a_j}$, which is not a constant phase.
At $m=2$, compare all-zero output with output $0100$, each against its
bitwise-complement input. This also excludes an identification merely by
periodic virtual gauge transformation or a global scalar.

The left-to-right pair $(a,b)$ is encoded as $2a+b$; successive pairs
give the four-bit order used for two neighboring blocked sites.\footnote{The
    existing encodings are \leanid{MPOTensor.CZX.siteBits} and
    \leanid{MPOTensor.CZX.localBits}.}
The formulas (\ref{eq:czx_repaired_block}) and
(\ref{eq:czx_decorated_block}) are an informal source audit, not new
public tensor definitions or a formal proof of their correspondence.

\section{Checked boundary and elimination plan}
The checked witness is exactly the printed tensor
(\ref{eq:czx_printed_tensor}). The formal results establish its two-site
all-zero row formula and its failure to generate a matrix product
unitary.\footnote{Witness: \leanid{MPOTensor.CZX.printed2017Tensor}.
    Results: \leanid{MPOTensor.CZX.printed2017Tensor_two_zero_row} and
    \leanid{MPOTensor.CZX.printed2017Tensor_not_isMPU}.}
The all-length contraction (\ref{eq:czx_periodic_contraction}) and the
comparison (\ref{eq:czx_decorated_block}) are mathematical derivations
in this note, not additional checked correspondence theorems.

The literal printed cross-paper tensor identity is false. The corrected
target is the exact Z-decorated tensor
(\ref{eq:czx_decorated_block}), with its all-length monomial formula and
unitarity, using ordinary blocking and the stated bit ordering.\footnote{The
    corrected target is tracked by \ghissue{7738} and its implementation
    child \ghissue{7752}; this audit does not assert that implementation
    is merged or that its correspondence theorem is available here.}
Any connection to the 2017 gate pair
$u=CZ(X\otimes X), v=CZ$, where $CZ=\operatorname{diag}(1,1,1,-1)$,
of \cite{Cirac2017MPU} must explicitly specify the normalization and
physical change of representative, rather than assert literal equality.
No classification signs, action data, defect tensors, or completion
certificates are supplied here. This audit does not establish the two
original CZX classification claims.
\section{The 2020 review repeats the printed tensor}
Appendix A, ``The CZX MPU'', of \cite{Cirac2021Matrix} writes the CZX
matrix product unitary as
$\lvert 0\rangle\langle 1\rvert\otimes|0)(+| +
    \lvert 1\rangle\langle 0\rvert\otimes|1)(-|$, with the normalized
vectors $(\pm|=2^{-1/2}\bigl((0|\pm(1|\bigr)$ on the virtual level. Its two
nonzero matrices are $2^{-1/2}A^{01}$ and $2^{-1/2}A^{10}$, where
\begin{align}
    A^{01} & =\begin{pmatrix}1&1\\0&0\end{pmatrix}, &
    A^{10} & =\begin{pmatrix}0&0\\1&-1\end{pmatrix}
    \label{eq:czx_review_matrices}
\end{align}
are the matrices the review prints two displays later. The bra-ket display
is exactly the tensor (\ref{eq:czx_printed_tensor}), and the printed
matrices (\ref{eq:czx_review_matrices}) are its rescaling
$A=\sqrt{2}P$.\footnote{Source traceability:
    \path{Papers/2011.12127/TN-Review-main.tex}, lines 2582--2601.
    Witnesses: \leanid{CZXCompression.reviewCZXTensor} and
    \leanid{CZXCompression.printed2017Tensor_eq_smul_reviewCZXTensor}.}

Every later claim of the passage holds for the matrices
(\ref{eq:czx_review_matrices}): the operator $O(A)$ is unitary and normal
after blocking two sites, its square has the displayed bond-four matrices
with integer entries, the invariant subspace of the projector $P$ and the
second reduction step with the projector $Q$ lead to the
canonical form $B^{ij}=(-1)\delta_{ij}$, and
$O(A)^2=(-1)^NI$.\footnote{Results:
    \leanid{CZXCompression.reviewCZXTensor_isMPUPos},
    \leanid{CZXCompression.reviewCZXMPS_isNormal},
    \leanid{CZXCompression.reviewCZXSquare_apply_zero},
    \leanid{CZXCompression.reviewQ_conj_reviewBlock} and
    \leanid{CZXCompression.mpo_reviewCZXTensor_mul_self}.}
For the normalized tensor of the bra-ket display the periodic operator is
$2^{-N/2}O(A)$, so its square is $2^{-N}(-1)^NI$ and the global identity
fails by the factor $2^{-N}$, as in (\ref{eq:czx_periodic_contraction}).
The formalization takes the printed matrices (\ref{eq:czx_review_matrices})
as the source tensor of the review and records the bra-ket display as this
normalization slip.

A smaller discrepancy of the same passage is recorded here.

\emph{Gate order.} The review describes the operator in words as
controlled-$Z$ gates between all adjacent sites ``followed by a Pauli $X$ on
all sites'', which in time order is $X^{\otimes N}D_N$, where $D_N$ is the
diagonal operator with entries $(-1)^{\sum_n s_ns_{n+1}}$ on a ring of $N$
sites; this is also the order $U_{CZX}=U_XU_{CZ}$ of \cite{ChenLiuWen2011CZX}.
The printed matrices (\ref{eq:czx_review_matrices}) contract instead to
$O(A)=D_NX^{\otimes N}$. Since $D_NX^{\otimes N}=(-1)^NX^{\otimes N}D_N$,
the two readings differ by the global sign $(-1)^N$. Both square to
$(-1)^NI$, so none of the later claims of the passage depends on the
choice; the formalization follows the printed matrices.\footnote{Source
    traceability: \path{Papers/2011.12127/TN-Review-main.tex}, lines
    2586--2591; \path{References/1106.4752/source/dDSPTmodel.tex},
    line 287. Results: \leanid{CZXCompression.mpo_reviewCZXTensor_apply} and
    \leanid{CZXCompression.mpo_reviewCZXTensor_eq_smul}.}

\bibliographystyle{alpha}
\bibliography{references}
\end{document}
